\begin{abstract}
Probabilistic Risk Assessment (PRA) is essential for risk-informed decision-making in industries such as nuclear, aerospace, oil and gas, chemical processing plants etc. However, the software tools used to analyze PRA models such as event trees, fault trees, and other risk models are often tied to specific operating systems and are hard to scale across different computing resources. This research work presents a domain specific distributed networked system that enables high throughput PRA quantification as a service. It provides a uniform, solver agnostic interface that brings multiple PRA engines together behind a single API gateway. The gateway isolates users from tool and operating system constraints and runs analyses on CPUs and GPUs across distributed memory and on-premise cluster. The research contributions can be divided into three parts. First, it provides a deployable microservice that runs PRA jobs asynchronously for multiple users and supports large volumes of jobs. It uses open source container orchestration to automate deployment, scaling, and routine management. Second, it includes an interoperability layer for PRA solvers that uses a common job specification and converts user submitted models to the supported solver formats automatically. Third, it introduces a three-step preprocessing, processing, and postprocessing pipeline. A case study of MHTGR models using the system shows that the pipeline speeds up analyses of large models while still producing reliable quantification results. The platform also provides controlled environments with role based access for regulated and proprietary data. For the nuclear engineering community, the system enables repeatable, high volume PRA quantification that are easier to run, compare, and share. 
%It also provides a software foundation that can add future solvers without disrupting existing analysis practices.
\end{abstract}

\thesistableofcontents
\thesislistoftables
\thesislistoffigures
